\section{From Exposure Adjustment to Returns}
\label{sec:identification-estimation}

\subsection{Scalar projection and the return equation}
\label{sec:exposures-to-returns}

The financial test asks whether return dependence follows the measured
information field.
Connecting that test to the exposure model requires a bridge from exposure
vectors to scalar returns.
Scalar projection carries the spatial operator and coefficient to the
systematic return component; a working quasi-likelihood then specifies the
observable return model.
The exposure vectors $\xi_i$ and $B_i$ and transmission map $T$ remain latent.
The application observes text laws $\widehat C_i$ and returns $r_t$, with
$W^\flat$ constructed from the pre-evaluation news window.
Projection onto a scalar risk direction preserves the exposure equilibrium.

\begin{proposition}[Scalar projection]\label{prop:scalar-projection}
  For every $h\in\mathbb R^K$, the scalar field $B_i^{\prime}h$ satisfies the
  ordinary spatial autoregression with the same $W$ and $\rho$, with reduced
  form $(1-\rho)(I-\rho W)^{-1}$ applied to $\xi_i^{\prime}h$.
\end{proposition}
By linearity of the dot product,
\[
  B_i^{\prime}h
  =\rho\sum_j W_{ij}B_j^{\prime}h
  +(1-\rho)\xi_i^{\prime}h.
\]
Stacking preserves the same $W$ and $\rho$; applying the inverse in
\Cref{thm:sar-reduced-form} gives the stated reduced form.

Define the systematic return component and the projected stand-alone exposure
by $s_t^i=B_i^{\prime}F_t$ and $z_t^i=\xi_i^{\prime}F_t$, and stack them
across firms.
Conditional on the realized factor direction $F_t$, scalar projection gives
$s_t=\rho Ws_t+(1-\rho)z_t$.
Since centered excess returns satisfy $\widetilde r_t=s_t+e_t$, substitution
yields

\begin{equation}\label{eq:model-return-innovation}
  \widetilde r_t
  =\rho W\widetilde r_t+x_t^{\mathrm{model}},
  \qquad
  x_t^{\mathrm{model}}
  =(1-\rho)z_t+(I-\rho W)e_t.
\end{equation}

Equation~\eqref{eq:model-return-innovation} is the return equation implied by
the exposure model.
Filtering the idiosyncratic components by $I-\rho W$ and adding the projected
stand-alone exposure generally makes $x_t^{\mathrm{model}}$ nonspherical and
correlated across firms, even when the components of $e_t$ are
cross-sectionally uncorrelated.
Allowing an intercept to absorb centering and mean components gives the
empirical notation

\begin{equation}\label{eq:spatial-lag}
  r_t = \rho W r_t + x_t .
\end{equation}

\subsection{Estimand}
\label{sec:estimand}

Three coefficients remain distinct throughout the paper.
The structural $\rho$ governs peer adjustment of latent exposures in
\Cref{sec:spatial-closure} and lies on the nonnegative-intensity branch $0\leq\rho<1$ of
\Cref{prop:peer-adjustment}, where $\lambda=\rho/(1-\rho)$.
For the concentrated working log-likelihood $\ell_t^Q(\rho;W)$ of
\Cref{sec:estimation}, the population target is the pseudo-true coefficient

\[
  \rho^\star
  \in\operatorname*{argmax}_{|\rho|<1}
  \E[\ell_t^Q(\rho;W)],
\]

and $\hat\rho$ is its sample estimate for a supplied $W$.
Identifying the population target with the structural coefficient requires the
population quasi-score to have its unique zero at the structural value despite
the richer innovation in \eqref{eq:model-return-innovation}.
This condition is additional to the algebraic result supplied by scalar
projection.

The fitted $\hat\rho$ measures conditional return dependence along the
supplied information structure.
Its transformation $\hat\lambda=\hat\rho/(1-\hat\rho)$ is an adjustment index
for this working return model.
Evaluation returns do not determine their own peer weights.
Text and returns can nevertheless share common determinants, so causal
interpretation requires an additional exclusion restriction or shock design.
This estimand and interpretation govern the empirical comparisons below.

\subsection{What the spatial field adds to a factor model}
\label{sec:factor-invariance}

The primary specification conditions excess returns on common factors with
unrestricted firm coefficients,
\begin{equation}\label{eq:factor-spatial-lag}
  (I-\rho W)r_t^e
  =\alpha+A f_t+\varepsilon_t,
  \qquad \varepsilon_t\mathrel{\overset{\mathrm{work}}{\sim}}
  \mathrm{iid}(0,\sigma^2I),
\end{equation}
where $r_t^e$ subtracts the daily risk-free rate, $\alpha$ contains
firm-specific intercepts, each row of $A$ contains one firm's loadings, $f_t$
contains the factor returns, and $W$ is the supplied interaction field,
$W^\flat$ in the primary specification.
Unrestricted firm coefficients pin down what the field can contribute.

\begin{proposition}[Factor-beta invariance]\label{prop:factor-beta-invariance}
  On a balanced panel with fixed $W$, common regressors $[1,f_t]$, and
  unrestricted firm coefficients, let $C_{\mathrm{OLS}}$ collect the
  coefficients of firm-by-firm least-squares regressions of $r_t^e$ on
  $[1,f_t]$.
  For every $\rho$ with $I-\rho W$ invertible, the profiled coefficients are
  $C(\rho)=(I-\rho W)C_{\mathrm{OLS}}$, and
  \begin{equation}\label{eq:factor-reduced-form-identity}
    (I-\rho W)^{-1}C(\rho)=C_{\mathrm{OLS}}.
  \end{equation}
\end{proposition}
Profiling applies least squares to $(I-\rho W)r_t^e$ on the common design, and
least squares is linear in the dependent variable, which gives $C(\rho)$ and
hence \eqref{eq:factor-reduced-form-identity}.

Under these conditions, standard factor exposures and fitted means remain
equal to their firm-by-firm regression estimates.
The field adds information through residual covariance: which firms'
unexpected returns move together, and how strongly.
The held-out comparisons evaluate how well alternative peer definitions
organize that joint distribution.
